\GalSourcePageBoundary{084}{L0083}{55}{55}
et de là il est aisé de conclure immédiatement que les substitutions
du groupe~$G$ doivent être toutes comprises dans la formule
\[
\tag{A}
\GalPrintedError{(a_{k_{1}, k_{2}},\ a_{m_{1} k_{1} + n_{1} k + \alpha_{1} m_{2} k_{1} + n_{2} k_{2} + \alpha_{2}})}{W10-PE005},
\]
Or nous savons, par le n\textsuperscript{o}~\GalHistoricalFootnoteMark{*},
que les substitutions du groupe~$G$
ne peuvent embrasser que $p^{2} - 1$ ou $p^{2} - p$ lettres. Ce n'est
point $p^{2} - p$, puisque, dans ce cas, le groupe~$G$ serait non primitif.
Si donc, dans le groupe~$G$, on ne considère que les permutations
où la lettre~$a_{0,0}$, par exemple, conserve toujours la même
place, on n'aura que des substitutions de l'ordre $p^{2} - 1$ entre les
$p^{2} - 1$ autres lettres.
\GalHistoricalFootnoteText{*}{Ce Mémoire faisant suite à un travail de Galois que je ne possède pas, il
m'est impossible d'indiquer le Mémoire cité ici et plus bas. \textup{(A.~Ch.)}}

Mais rappelons-nous ici que c'est simplement pour la démonstration
que nous avons supposé que le groupe primitif~$G$ se partageât
en groupes conjugués non primitifs. Comme cette condition
n'est nullement nécessaire, les groupes seront souvent beaucoup
plus composés.

Il s'agit donc de reconnaître dans quel cas ces groupes pourront
admettre des substitutions où $p^{2} - p$ lettres seulement varieraient,
et cette recherche va nous retenir quelque temps.

Soit donc $G$ un groupe qui contienne quelque substitution de
l'ordre $p^{2} - p$; je dis d'abord que toutes les substitutions de ce
groupe seront linéaires, c'est-à-dire de la forme~\textup{(A)}.

La chose est reconnue vraie pour les substitutions de l'ordre
$p^{2} - 1$; il suffit donc de la démontrer pour celles de l'ordre
$p^{2} - p$. Ne considérons donc qu'un groupe où les substitutions
seraient toutes $m$ de l'ordre~$p^{2}$ ou de l'ordre $p^{2} - p$. (\emph{Voyez} l'endroit
cité.)

Alors les $p$~lettres qui, dans une substitution de l'ordre $p^{2} - p$,
ne varieront pas, devront être des lettres conjointes.

Supposons que ces lettres conjointes soient
\[
a_{0,0},\quad a_{0,1},\quad a_{0,2},\quad \dots, \quad a_{0,p-1}.
\]
Nous pouvons déduire toutes les substitutions où ces $p$~lettres ne
changent pas de place, nous pouvons les déduire de substitutions
